
%% Revisited JRSI manuscript in English, adapted to the provided template
\documentclass[openacc]{rsproca_new}

\usepackage[T1]{fontenc}
\usepackage[utf8]{inputenc}
\usepackage{lmodern}
\usepackage{microtype}
\usepackage{graphicx}
\usepackage{booktabs}
\usepackage{longtable}
\usepackage{array}
\usepackage{amsmath,amssymb}
\usepackage{threeparttable}
\usepackage{multicol}
\usepackage{url}
\usepackage{csquotes}

\addbibresource{refs.bib}

\jname{rsif}
\Journal{J R Soc Interface\ }

\begin{document}

\title{Dynamic compliance in a balanced pelvic ring: sex-specific organisation of childbirth-relevant deformation pathways}

\author{Petr Heny\v{s}$^{1}$, Michal Kucha\v{r}$^{1,2}$, Alexander Mor\'{a}vek$^{2}$, Niels Hammer$^{3,4}$}

\address{$^{1}$Institute of New Technologies and Applied Informatics, Faculty of Mechatronics, Informatics and Interdisciplinary Studies, Technical University of Liberec, Liberec, Czech Republic\\
$^{2}$Department of Anatomy, Faculty of Medicine in Hradec Kr\'{a}lov\'{e}, Charles University, Hradec Kr\'{a}lov\'{e}, Czech Republic\\
$^{3}$Division of Macroscopic and Clinical Anatomy, Gottfried Schatz Research Center, Medical University of Graz, Graz, Austria\\
$^{4}$Division of Macroscopic and Clinical Anatomy, Gottfried Schatz Research Center, Medical University of Graz, Graz, Austria; Department of Orthopedic and Trauma Surgery, University of Leipzig, Leipzig, Germany; Division of Biomechatronics, Fraunhofer IWU, Dresden, Germany}

\subject{biomechanics, evolutionary anthropology, obstetrics}

\keywords{biomechanics, human pelvis, modal analysis, finite-element modelling, childbirth, allometry, sacroiliac joint}

\corres{Michal Kucha\v{r}\\
\email{kucharm@lfhk.cuni.cz}}

\begin{abstract}
The human pelvis must simultaneously support habitual bipedal loading and accommodate the rare but extreme mechanical demands of childbirth. This dual role is usually discussed in terms of sexual dimorphism and static canal dimensions, yet the pelvis functions mechanically as a compliant ring in which load transfer depends on the sacroiliac joints, the pubic symphysis and the accessibility of specific deformation pathways. We analysed 278 anatomically standardised lumbopelvic finite-element models derived from clinical CT using correspondence-preserving modal decomposition, allometric normalisation, standardised locomotor and childbirth-motivated load cases, and sensitivity analyses of ligamentous and cartilaginous constraints. Pelvic size explained a substantial share of global stiffness variation, indicating that size is a major determinant of pathway accessibility. After adjustment for scale and age, however, sex differences persisted in selected modes, most prominently modes~2 and~9, showing that sex reorganises ring mechanics rather than merely scaling them. Under bilateral and unilateral stance, both sexes recruited a shared low-order support backbone. Under simulated inlet opening, mode~9 emerged as the principal inlet-opening pathway, but men concentrated a larger fraction of the response in this single pathway whereas women distributed deformation more broadly across coordinated low-order modes. Sensitivity analyses localised regulation of mode~9 to the pubic symphysis, anterior sacroiliac ligaments and sacroiliac cartilage interfaces. We therefore interpret the female pelvis not as generally lax or evolutionarily compromised, but as a highly specified balanced compliant ring that preserves habitual support while retaining selective access to childbirth-relevant deformation during parturition.
\end{abstract}

\maketitle

\begin{multicols}{2}

\section{Introduction}

The human pelvis sits at the intersection of biomechanics, evolutionary anthropology and obstetrics because it must satisfy two apparently incompatible demands. It must transmit loads continuously between the trunk and lower limbs during stance and locomotion, yet it must also define the bony boundaries of a birth canal through which a large-brained neonate must pass. For more than a century, that dual role has encouraged a predominantly geometric interpretation of pelvic function. In the classical telling, childbirth is difficult because selection for efficient bipedalism constrains the pelvis, whereas selection for parturition favours widening of the canal \cite{RosenbergTrevathan2002,Trevathan2015,Grunstra2023}. In this view, pelvic dimensions themselves become the decisive currency of function.

That framing remains useful, but it is incomplete. The pelvis is not merely a set of distances. It is a mechanically coupled ring in which load transfer depends on the interplay between geometry, material distribution, cartilage interfaces and ligamentous constraints. A pelvis does not function as a static assemblage of diameters in the clinic, in locomotion or during parturition. It functions as a deformable structure whose behaviour depends on which deformation pathways are available, how strongly those pathways are recruited by a given load, and which anatomical constraints regulate access to them.

Several strands of literature make this broader perspective necessary. Experimental work has shown that wider female pelves do not inevitably increase locomotor cost \cite{Warrener2015}. At the same time, recent reviews have made clear that the obstetrical dilemma should neither be discarded as a myth nor reduced to a one-dimensional opposition between locomotion and parturition \cite{Pavlicev2020,Haeusler2021,Grunstra2023}. Instead, the pelvis must satisfy several partially competing demands at once, including load transfer, childbirth, pelvic-floor support and developmental robustness \cite{Grunstra2019,Haeusler2021}. Population studies have further shown that the shape of the birth canal varies substantially across and within human populations \cite{BettiManica2018,Auerbach2018}, while ontogenetic and developmental work indicates that adult pelvic form remains historically and biologically contingent rather than mechanically predetermined by one selective pressure alone \cite{VerbruggenNowlan2017,Novak2020,KurkiDecrausaz2015}. Taken together, these observations imply that static size alone cannot account for pelvic function.

What is still missing is a framework that captures pelvic mechanics directly at population scale. Subject-specific finite-element models have already demonstrated that sacroiliac constraints, pubic symphysis stiffness and load direction materially affect pelvic behaviour \cite{Eichenseer2011,Hammer2013}. This is biomechanically plausible because the sacroiliac joint is a specialised load-transfer articulation with limited but functionally meaningful motion, and because its ligamentous architecture is strongly involved in force closure and stability \cite{Vleeming2012,Kiapour2019,ChoKwak2020}. Yet isolated subject-specific simulations do not easily reveal which aspects of mechanical behaviour are shared across individuals, which vary with scale, and which are structured by sex. If the pelvis is to be analysed as a biological system rather than as a collection of separate engineering case studies, we need a way to preserve anatomical correspondence across individuals while describing each pelvis through its intrinsic deformation repertoire.

Modal analysis provides such a framework. Instead of reducing the pelvis to a few linear dimensions, modal decomposition characterises it through its preferred deformation modes, each associated with a stiffness cost. Static or quasi-static loading scenarios can then be projected into the modal basis, allowing pelvic function to be described as the recruitment of deformation pathways rather than the displacement of a few landmarks. This shift changes the biological question. Instead of asking whether a pelvis is merely wide or narrow, we may ask which pathways are mechanically accessible, how they are organised, and whether obstetrically relevant deformation depends on the same modes that govern habitual support mechanics.

Our analysis is organised around a single central idea. We propose that the female pelvis is best understood not as an evolutionarily compromised or generally loosened structure, but as a highly specified \emph{balanced compliant ring}. In this view, habitual support remains built upon a largely shared low-order backbone in both sexes, while childbirth-relevant compliance is achieved through selective access to particular deformation pathways. Pelvic size is expected to modulate the accessibility of those pathways, whereas sex is expected to modulate how the ring organises response once those pathways are engaged. Those pathways, in turn, are expected to be regulated not by gross widening alone but by an anterior constraint architecture involving the pubic symphysis and the sacroiliac interfaces.

We therefore test four linked propositions. First, pelvic size should explain a substantial fraction of stiffness variation and thereby organise the baseline accessibility landscape, but not all of it. Second, some sex differences should remain after adjustment for scale and age, indicating a genuine reorganisation of ring mechanics rather than a trivial allometric effect. Third, locomotor load cases should recruit a shared low-order modal support basis, whereas childbirth-motivated load cases should reveal sex-biased organisation of response within that basis. Fourth, the principal inlet-opening pathway should be governed primarily by anterior constraints rather than by static dimensions alone. If supported, these propositions would imply that pelvic function is best described not by geometry in isolation but by the architecture of dynamic compliance.

\section{Material and methods}

\subsection{Cohort and correspondence-preserving framework}

The analysis used the BoneDat resource, a standardised database of 278 clinical lumbopelvic CT scans from a Central European cohort spanning adolescence to old age \cite{HenysKuchar2025}. In the present derived workbooks, the cohort comprised 150 female and 128 male individuals, aged 16--91 years. BoneDat provides segmented anatomy, a shape-normalisation pipeline, a common reference surface and volumetric meshes, allowing each subject to be represented on a shared anatomical domain while preserving subject-specific morphology through deformation fields. This correspondence-preserving design is essential here because it permits direct mode-wise and region-wise comparison across the full cohort.

The original CT data were acquired on a Siemens Definition AS+ scanner and passed the technical quality control reported for BoneDat. Because the scans were acquired without a calibration phantom, material properties should be interpreted as cohort-wise relative descriptors rather than direct clinical densitometry. Phantom-less internal calibration was used to convert the released HU values to hydroxyapatite-equivalent density before elastic property assignment.

\subsection{Allometry and static morphology}

Allometric descriptors were taken from \texttt{allometry.xlsx}. The principal size variable was defined as the cube root of total pelvic volume relative to a template reference volume, which provides an appropriate whole-pelvis scale measure for a three-dimensional structure. Additional descriptors included total surface area, total mass, effective thickness and effective density. Static morphological dimensions were extracted from \texttt{anatomy\_data.xlsx}; the main variables were the anteroposterior inlet diameter (AP), pelvic inlet transverse diameter (PIT), subpubic angle, biacetabular width, biischiadic width, bituberous width, sacral width and iliopectineal eminence width. These variables were retained as the conventional pelvimetric descriptors against which the mechanical phenotype could be compared.

\subsection{Finite-element modelling and modal decomposition}

Each released deformation field was used to warp a tetrahedral template mesh into subject space, after which mechanics were evaluated on the common reference domain. Bone, sacroiliac cartilage and the pubic symphysis were represented on this shared mesh. Bone material properties were spatially heterogeneous and derived from mapped hydroxyapatite-equivalent density; sacroiliac cartilage and the pubic symphysis were assigned region-wise material properties. Linear isotropic elasticity was used throughout.

For each subject, pelvic mechanics were summarised through a generalised elastic eigenproblem,
\begin{equation}
K\phi_{k} = \lambda_{k} M\phi_{k}, \qquad k=1,2,\ldots,
\end{equation}
where $K$ and $M$ are the subject-specific stiffness and mass matrices on the mapped domain. Lower eigenvalues were interpreted as mechanically more accessible deformation pathways. We focused on modes~1--10, which the supplied documentation identifies as the biologically relevant low-order spectrum.

This modal framework is central to the present interpretation because it separates two distinct features of pelvic behaviour. \emph{Pathway accessibility} refers to where a mode lies in the stiffness spectrum: low-order, low-eigenvalue modes are easier to recruit mechanically. \emph{Pathway dominance} refers instead to how strongly a given load case concentrates response energy into that mode once the pathway has been engaged. This distinction is especially important for mode~9: it is interpreted here as the principal inlet-opening pathway, but not necessarily as the energetically most dominant component in every subject or sex.

\subsection{Load cases and modal participation}

Static responses were computed for five standardised loading scenarios. Two represented habitual support mechanics: bilateral stance (\texttt{SP2leg}) and unilateral stance (\texttt{SP1leg}). Three represented childbirth-motivated probes: inlet opening (\texttt{LAB\_phase1}), outlet opening (\texttt{LAB\_phase2}) and terminal descent (\texttt{LAB\_phase3}). These scenarios were not intended as literal reproductions of labour, but as controlled mechanical probes that distinguish support-related from childbirth-relevant deformation pathways.

Static displacement fields were projected onto the modal basis by solving
\begin{equation}
(\Phi^{T}M\Phi)q = \Phi^{T}Mu,
\end{equation}
where $u$ is the static response, $\Phi=[\phi_{1},\ldots,\phi_{m}]$ is the modal basis, and $q$ is the vector of modal participation factors. Modal strain energies were then computed as
\begin{equation}
U_{i}=\frac{1}{2}\lambda_{i}q_{i}^{2}\|\phi_{i}\|_{M}^{2}, \qquad
\eta_{i}=\frac{U_{i}}{\sum_{j}U_{j}}.
\end{equation}
For interpretation and table construction, fractions were normalised over modes~1--10 so that the load-case summaries remained focused on the mechanically meaningful low-order spectrum.

\subsection{Sensitivity analysis}

To identify which anatomical constraints regulate the most relevant deformation pathways, we analysed workbook-based sensitivities of eigenvalues to ligament stiffness and cartilage-related modulus changes. For a ligament bundle with axial stiffness $EA_{\ell}$, the first-order modal sensitivity was
\begin{equation}
\frac{\partial \lambda_{i}}{\partial EA_{\ell}}
=
\frac{\phi_{i}^{T}(\partial K/\partial EA_{\ell})\phi_{i}}
{\phi_{i}^{T}M\phi_{i}}.
\end{equation}
For cartilage and symphyseal regions, the code reported a normalised energy-based sensitivity proxy,
\begin{equation}
S_{i,R}=\frac{\mathcal{E}_{i}(R)}{E_{R}\,\phi_{i}^{T}M\phi_{i}},
\end{equation}
which is proportional under linearisation to $\partial \lambda_{i}/\partial E_{R}$. These metrics were used to rank the importance of ligament groups, sacroiliac cartilage interfaces and the pubic symphysis for individual modes. Our primary mechanistic focus was mode~9, with modes~2 and~5 retained as support-related and outlet-related reference modes, respectively.

\subsection{Statistical analysis}

Static sex differences were assessed with Welch two-sample $t$-tests. Mode-wise sex effects were estimated using linear models of the form
\begin{equation}
\log(\lambda) \sim \log(\mathrm{scale}) + \mathrm{sex} + \mathrm{age}.
\end{equation}
The corresponding sex coefficient was interpreted as the percentage change in modal stiffness for males relative to females after adjustment for scale and age. To evaluate the limits of static pelvimetry, we also fitted female-only univariate models relating size-corrected eigenvalues of modes~9 and~2 to individual pelvic dimensions.

\section{Results}


\subsection{Cohort structure and static dimorphism}

The cohort comprised 278 individuals, including 150 women and 128 men, with no statistically meaningful age difference between the sexes (table~\ref{tab:cohort}). As expected, female pelves were smaller in overall scale than male pelves ($0.925 \pm 0.045$ versus $0.998 \pm 0.045$, $p=6.51 \times 10^{-32}$), yet they showed the classic enlargement of childbirth-relevant dimensions. AP and PIT were both larger in women, the subpubic angle was markedly wider, and both biischiadic and bituberous widths were substantially greater. These differences confirm that the dataset captures well-known sexual dimorphism in static pelvic morphology.

At the same time, not all regions behaved uniformly. Sacral width showed no meaningful sex difference in the harmonised workbook values, whereas several anterior and outlet-related dimensions did. This is already suggestive. Pelvic dimorphism does not appear to be a simple global expansion of the female pelvis, but a regional reorganisation whose functional meaning must be read mechanically.

\subsection{Allometry explains much, but not everything}

Allometric regressions showed that total volume scaled essentially isometrically with scale (slope $=3.000$, $R^2=1.000$), whereas total surface area scaled with a shallower exponent (slope $=1.684$, $R^2=0.894$). Total mass increased with scale (slope $=2.679$, $R^2=0.902$), while effective density decreased moderately with scale (slope $=-0.321$, $R^2=0.117$). In other words, larger pelves were not simply magnified replicas of smaller pelves. Their geometric and material characteristics changed in a way that implies non-trivial mechanical consequences.

This matters because size is a major determinant of stiffness. A smaller ring will generally be more compliant than a larger one, all else being equal. If sex differences in pelvic mechanics were due solely to this size effect, then scale correction should eliminate them.

\subsection{Sex differences persist after size correction}

That did not occur. Several modes retained clear sex differences after adjustment for scale and age (table~\ref{tab:modes}). The largest adjusted differences occurred in mode~2 and mode~9, where male pelves were stiffer than female pelves by $20.72\%$ and $20.07\%$, respectively. Additional significant sex effects were present in modes~3, 4, 7, 8 and~10. Mode~1 remained effectively neutral.

This pattern is central to the biological interpretation. Female compliance is not merely a by-product of smaller average size. Instead, the stiffness spectrum is redistributed in a mode-specific way. Some deformation pathways remain intrinsically more accessible in female pelves even after scale is removed from the comparison.

Mode~1 is especially informative because it serves as a near-neutral reference. Its lack of a meaningful sex effect suggests that the female pelvis is not globally re-engineered for childbirth. Rather, only selected pathways differ. Modes~2 and~9, by contrast, indicate that sex-specific reorganisation affects both baseline ring support and childbirth-relevant compliance.

\subsection{Static pelvimetry explains only a modest fraction of modal behaviour}

We next asked whether the two most relevant modes could be recovered from simple pelvic dimensions. In women, PIT explained $13.5\%$ of the variation in size-scaled mode~9 eigenvalues, while AP explained $9.2\%$. For mode~2, the strongest single predictor was PIT ($R^2=0.089$), followed by bituberous width, biischiadic width and biacetabular width, all of which explained less than $8\%$ of the variance (table~\ref{tab:staticassoc}).

No single static dimension explained even one fifth of the relevant mechanical variability. This does not make pelvimetry useless. PIT and AP clearly retain information. But it does show that pelvic mechanics cannot be inferred from any single diameter. Mechanical phenotype is distributed across geometry, material mapping and constraint structure.

\subsection{Locomotor load cases rely on a shared low-order support backbone}

Under bilateral stance, both sexes recruited an overwhelmingly similar response dominated by mode~2: $84.6\%$ in women and $84.5\%$ in men (table~\ref{tab:loadcases}). Secondary participation was minor. This is the clearest expression of a shared support-related ring mode.

Unilateral stance broadened the response but preserved the same logic. Both sexes relied primarily on modes~1--3. In women, the mean participation fractions were $39.6\%$ for mode~1, $37.8\%$ for mode~2 and $15.3\%$ for mode~3; in men the corresponding values were $40.4\%$, $37.2\%$ and $14.1\%$. Again, sex differences were small.

These results imply that habitual support mechanics are built on a largely conserved low-order modal architecture. Everyday load transfer does not require fundamentally different deformation strategies in women and men.

\subsection{Inlet opening reveals sex-specific organisation of compliance}

The picture changed sharply under the inlet-opening scenario. Mode~9 was the dominant response component in 265 of the 278 subjects and is therefore the principal inlet-opening pathway in the present cohort and modelling framework. However, its contribution differed between sexes in a way that is easy to misread unless pathway accessibility and energy concentration are kept separate.

Men placed a larger mean fraction of the response into mode~9 ($52.5\%$) than women ($46.4\%$). Women, however, recruited much more of mode~4 ($18.2\%$ versus $6.0\%$) and maintained a broader spread of low-order participation. Men, by contrast, showed relatively greater concentration into mode~9 and somewhat greater high-order involvement through modes~8 and~10.

The important point is therefore not that women have a stronger mode~9 in fractional terms. They do not. Rather, mode~9 is the principal inlet-opening pathway, whereas women appear to organise inlet-opening behaviour more distributively across a coordinated low-order set once that pathway class is engaged. Men reach a similar class of deformation with a greater concentration into one dominant pathway. This distinction underlies our interpretation of the female pelvis as a balanced compliant ring rather than as a generally more lax structure.

\subsection{Outlet opening is organised differently but remains linked to inlet-related pathways}

Under outlet-opening loading, the dominant mode was mode~5 in both sexes ($40.7\%$ in women, $39.2\%$ in men), identifying it as the best candidate for an outlet-related deformation pathway. Yet this response was not independent of inlet mechanics. Women retained a greater concurrent contribution of mode~9 ($16.1\%$ versus $12.4\%$), whereas men showed larger participation of mode~6.

This suggests that inlet- and outlet-related mechanics are not separate modules. The female pelvis appears to maintain a broader low-order compliance repertoire even when the dominant scenario shifts away from inlet opening.

\subsection{Terminal descent returns to the support-related backbone}

In the terminal descent scenario, mode~2 again dominated almost universally ($77.2\%$ in women, $74.7\%$ in men), with mode~8 acting as the principal secondary contributor. This reinforces the interpretation that support-related ring mechanics remain a central substrate even in childbirth-related scenarios. Obstetric function does not arise in a mechanically isolated system. It is layered onto the same ring architecture that governs everyday support.

\subsection{Mode~9 is governed by an anterior constraint system}

Sensitivity analysis provided the strongest mechanistic insight. At the interface level, mode~9 was most sensitive to the left sacroiliac cartilage interface, followed by the right sacroiliac cartilage interface and the pubic symphysis interface. At the ligament level, the largest effect was associated with the symphyseal ligament group, followed by the anterior sacroiliac ligaments and, at a lower magnitude, the interosseous sacroiliac ligaments. Posterior sacroiliac ligaments, sacrospinous ligaments and sacrotuberous ligaments were far less influential (table~\ref{tab:sensitivity}).

Mechanically, this means that the principal inlet-opening pathway is not determined primarily by gross bony width. It is gated by an anteriorly weighted constraint system spanning the pubic symphysis and the sacroiliac interfaces. The pathway therefore belongs to the mechanics of the whole ring, not merely to the geometry of the inlet.

\section{Discussion}

The principal contribution of this study is not that it reiterates well-known static dimorphism in pelvic dimensions. Those differences are real, and they are present in our data, but they are not the main discovery. The central result is that even after the removal of size effects, the stiffness spectrum of the human pelvis remains sexually structured in a mechanically interpretable and functionally meaningful way. The female pelvis is therefore not simply a smaller or wider version of the male pelvis. It is a differently organised ring.

That conclusion changes the level at which pelvic function should be discussed. Instead of treating function as a direct readout of canal dimensions, it becomes more productive to treat the pelvis as a compliant mechanical system with a structured repertoire of accessible deformations. This is why the language of the \emph{balanced compliant ring} is useful. It does not deny morphology. It situates morphology within a broader mechanical architecture.

The first necessary step toward this conclusion was allometric separation. Size is a major driver of pelvic stiffness, and any attempt to discuss female compliance without accounting for scale risks triviality. Smaller structures will generally deform more easily. Yet the persistence of strong adjusted sex effects in modes~2 and~9 shows that size alone is insufficient. There is a genuine reorganisation of how stiffness is distributed across deformation pathways.

Mode~9 is the clearest case. It is sex-biased but not female-exclusive. It is strongly recruited by simulated inlet opening but not by habitual stance. It is only modestly predicted by static inlet dimensions. And it is governed primarily by the pubic symphysis and anterior sacroiliac structures. That combination of features makes it more than a statistically interesting eigenmode. It makes it the best candidate for an inlet-opening pathway at the level of the whole pelvic ring.

It is important, however, not to replace one reduction with another. We do not argue that mode~9 is \emph{the} obstetric mode in any absolute sense. Childbirth is too complex, and the models here are too idealised, for such a claim. Rather, mode~9 is the dominant mechanically interpretable route by which inlet-opening behaviour appears in the present dataset. Its biological relevance comes not from exclusivity but from its position within a network of related low-order modes and its reproducible association with obstetric loading.

This is why the load-case results matter so much. Men and women did not differ primarily in whether mode~9 existed. They differed in how inlet-opening response was organised around it. Men concentrated a greater fraction of energy into mode~9, whereas women distributed the response more broadly across mode~9, mode~4 and additional low-order components. This is the key empirical basis for the notion of distributed compliance in the female pelvis.

That pattern should not be misinterpreted as diffuse laxity. On the contrary, diffuse laxity would be mechanically undesirable because it would threaten everyday support stability. Our findings point toward something much more selective. Women retain a locomotor backbone that is remarkably similar to that of men under bilateral and unilateral stance. What differs is that childbirth-relevant loading reveals access to a wider coordinated set of low-order pathways. In that sense, the female pelvis appears not softer in general, but \emph{selectively more available} in specific directions of deformation.

This interpretation helps move beyond the simplest version of the obstetrical-dilemma narrative without discarding biomechanics altogether. Experimental evidence suggests that wider pelves need not impose a major locomotor penalty \cite{Warrener2015}, and recent syntheses argue that childbirth must be understood as the product of multiple interacting pressures rather than a single fixed trade-off \cite{Pavlicev2020,Haeusler2021,Grunstra2023}. Our results fit this broader view. They indicate that locomotor stability and childbirth-relevant compliance need not lie on a single axis. Instead, they may be distributed across different layers of pelvic mechanical organisation: a shared support backbone and a selectively modulated compliance architecture. In that sense, our modal framework offers a mechanical language for ideas that have usually been expressed only morphologically or verbally: the same pelvis can remain stable in habitual support while preserving conditional access to opening-related deformation.

The sensitivity analysis strengthens this argument. If the principal inlet-opening pathway is most strongly governed by the pubic symphysis, anterior sacroiliac ligaments and sacroiliac cartilage interfaces, then childbirth-relevant compliance is best understood as the product of targeted tuning within the anterior ring. This is mechanically elegant because it allows the pelvis to preserve overall integrity while maintaining conditional access to opening-related deformation. Posterior structures can continue to stabilise the system, while the anterior constraint complex regulates the pathway most relevant to the inlet. This interpretation is also consistent with review literature describing the human pubic symphysis as robust, only modestly flexible under physiological conditions, yet still biologically meaningful for childbirth and pelvic stability \cite{Grunstra2019,Stolarczyk2021}. Our results refine that point by showing where such limited mobility is likely to matter most: not as diffuse ring laxity, but as selective access to a coordinated deformation pathway.

That idea resonates with previous finite-element work showing that ligamentous architecture materially affects pelvic load distribution \cite{Eichenseer2011,Hammer2013}, and with broader biomechanical reviews of SIJ function that emphasise small but consequential joint motion in a heavily constrained system \cite{Vleeming2012,Kiapour2019,ChoKwak2020}. It also offers a plausible biomechanical substrate for broader comparative and obstetric observations that flexibility and posture can modulate cephalopelvic constraint without eliminating it \cite{Pink2024,Fremondiere2022}. Our contribution is to show that such flexibility need not be invoked as a vague property of the whole pelvis. It can be grounded in a specific hierarchy of deformation pathways and constraints.

The limited explanatory power of static pelvimetry deserves emphasis. PIT and AP were the strongest single predictors of mode~9, but even they accounted for only modest fractions of the variance. This means that pelvimetry remains informative, yet mechanically incomplete. A pelvis may present apparently favourable dimensions while still differing substantially in the accessibility of key deformation pathways. For biomechanics, anthropology and potentially clinical interpretation, that is a major conceptual shift.

From an evolutionary perspective, the implications are equally important. If mechanical performance depends not only on canal size but on the organisation of accessible deformation pathways, then selection may have acted on the architecture of compliance itself. This permits a richer interpretation of human pelvic variation. Appreciable diversity of static form can coexist with functional robustness if different shapes still preserve a suitable low-order mechanical repertoire. That possibility is consistent with the documented geographic structuring of pelvic form \cite{BettiManica2018,Auerbach2018}, with broader arguments that obstetric adaptation is distributed across the pelvis rather than localised to a single dimension \cite{KurkiDecrausaz2015}, and with developmental influences on adult pelvic dimensions \cite{VerbruggenNowlan2017,Novak2020}.

Several limitations must be stated clearly. The finite-element models are linear and evaluate small-deformation behaviour. They do not capture full nonlinear labour mechanics, hormonal softening in pregnancy, viscoelasticity, maternal posture change, muscle activity or fetal interaction. The cohort is a morphology dataset rather than an obstetric outcomes cohort, and it represents a Central European clinical population rather than all human diversity. The precise ethical approval number, final funding statement and public DOI for code and derived tables must be inserted before submission. In addition, the variable documentation referenced a \texttt{pretension\_sensitivity.xlsx} workbook that was not included among the uploaded files, so the present interpretation is explicitly limited to stiffness and interface sensitivities.

These limitations define the next step rather than weaken the present result. What this study establishes is a coherent population-level mechanical framework in which pelvic function can be described through modal phenotype. The natural extension is toward outcome-linked, pregnancy-aware and possibly nonlinear studies that test whether the same compliance pathways matter in vivo. In particular, future work should connect the modal phenotype of the bony ring to pelvic-floor state and dynamic birth simulations, because the literature increasingly suggests that obstetric performance emerges from coupled skeletal and soft-tissue mechanics rather than from bone geometry alone \cite{Pavlicev2020,Fremondiere2022}.

\section{Conclusion}

In a population-scale, correspondence-preserving analysis of 278 human pelves, we found that size explains much of global stiffness but does not fully account for pelvic mechanics. After adjustment for scale and age, strong sex differences remain in selected deformation modes, especially modes~2 and~9. Habitual support scenarios rely on a shared low-order modal backbone in both sexes, whereas childbirth-motivated loading reveals sex-specific reorganisation of compliance. The principal inlet-opening pathway is governed primarily by the pubic symphysis and anterior sacroiliac constraint system.

We therefore argue that the female pelvis is best understood not as an evolutionarily compromised or generally lax structure, but as a highly specified balanced compliant ring: a structure that preserves everyday support and locomotor loading while maintaining selectively accessible, anteriorly regulated pathways for childbirth-relevant deformation. This reframes pelvic biomechanics from a question of static geometry to one of mechanical architecture and dynamic compliance.

\section*{Ethics}

This study used retrospectively acquired clinical CT data released through the BoneDat resource. According to the BoneDat dataset documentation, the study was approved by the Ethics Committee of University Hospital Hradec Kr'{a}lov'{e} (Ref: 202411I03P).

\section*{Data accessibility}

The source morphology dataset is described in BoneDat \cite{HenysKuchar2025}. Before submission, we recommend depositing derived statistical tables, figure-generation scripts and supplementary processing code in a public repository and inserting the repository DOI here.

\section*{Authors' contributions}

P.H. and M.K. designed the study. P.H. and M.K. performed the computational workflow. P.H., M.K. and A.M. analysed the data. P.H., M.K. and A.M. drafted the manuscript. N.H. contributed anatomical and biomechanical interpretation and critically revised the text. All authors approved the final version.

\section*{Competing interests}

The authors declare no competing interests.

\section*{Funding}

[Insert final funding statement and grant numbers.]

\section*{Acknowledgements}

We thank the BoneDat project for providing the standardised anatomical framework on which this analysis is based. Additional technical and institutional acknowledgements may be added as required for the final submission version.

\end{multicols}

\begin{table}[p]
\centering
\caption{Cohort summary and static pelvic measures. Values are mean $\pm$ s.d. $p$-values are Welch two-sample $t$-tests (female versus male). AP, anteroposterior inlet diameter; PIT, pelvic inlet transverse diameter.}
\label{tab:cohort}
\begin{tabular}{lccc}
\toprule
Measure & Female $(n=150)$ & Male $(n=128)$ & $p$ \\
\midrule
Age (years) & $53.4 \pm 19.7$ & $54.5 \pm 18.9$ & 0.617 \\
Scale & $0.925 \pm 0.045$ & $0.998 \pm 0.045$ & $6.51 \times 10^{-32}$ \\
Effective density & $1.341 \pm 0.079$ & $1.321 \pm 0.074$ & 0.028 \\
AP (mm) & $121.0 \pm 10.0$ & $112.3 \pm 9.1$ & $4.63 \times 10^{-13}$ \\
PIT (mm) & $134.9 \pm 8.6$ & $127.7 \pm 7.6$ & $2.88 \times 10^{-12}$ \\
Subpubic angle (deg) & $85.8 \pm 6.1$ & $69.1 \pm 5.3$ & $7.30 \times 10^{-71}$ \\
Biacetabular width (mm) & $129.1 \pm 8.2$ & $122.0 \pm 7.2$ & $2.94 \times 10^{-13}$ \\
Biischiadic width (mm) & $109.9 \pm 7.1$ & $97.7 \pm 7.1$ & $9.40 \times 10^{-35}$ \\
Bituberous width (mm) & $136.5 \pm 8.3$ & $127.3 \pm 7.5$ & $3.03 \times 10^{-19}$ \\
Sacral width (mm) & $116.7 \pm 5.7$ & $117.5 \pm 6.5$ & 0.327 \\
Iliopectineal eminence width (mm) & $106.2 \pm 7.0$ & $97.9 \pm 6.0$ & $1.90 \times 10^{-22}$ \\
\bottomrule
\end{tabular}
\end{table}

\begin{table}[p]
\centering
\caption{Allometric scaling of geometric and material-derived quantities. Slopes are from log--log regressions against pelvic scale.}
\label{tab:allometry}
\begin{tabular}{lccc}
\toprule
Outcome versus scale & Slope & $p$ & $R^2$ \\
\midrule
Total volume & 3.000 & $<10^{-300}$ & 1.000 \\
Total surface & 1.684 & $2.53 \times 10^{-136}$ & 0.894 \\
Total mass & 2.679 & $2.48 \times 10^{-140}$ & 0.902 \\
Effective density & $-0.321$ & $5.07 \times 10^{-9}$ & 0.117 \\
\bottomrule
\end{tabular}
\end{table}

\begin{table}[p]
\centering
\caption{Mode-wise scale dependence and adjusted sex effects (modes~1--10). Correlations are Pearson's $r$ between $\log(\mathrm{scale})$ and $\log(\lambda)$ for unscaled and size-scaled eigenvalues. Adjusted sex effect is percent change (male relative to female) from $\log(\lambda) \sim \log(\mathrm{scale}) + \mathrm{sex} + \mathrm{age}$.}
\label{tab:modes}
\begin{tabular}{cccccc}
\toprule
Mode & $r$ (unscaled) & $p$ & $r$ (scaled) & $p$ & Sex effect \\
\midrule
1 & $-0.461$ & $4.96 \times 10^{-16}$ & 0.009 & 0.876 & $+0.85\%$ ($p=0.734$) \\
2 & $-0.331$ & $1.52 \times 10^{-8}$ & 0.374 & $1.20 \times 10^{-10}$ & $+20.72\%$ ($p=2.62 \times 10^{-18}$) \\
3 & $-0.467$ & $1.82 \times 10^{-16}$ & 0.233 & $8.84 \times 10^{-5}$ & $+11.23\%$ ($p=1.80 \times 10^{-7}$) \\
4 & $-0.509$ & $1.04 \times 10^{-19}$ & 0.273 & $3.70 \times 10^{-6}$ & $+15.90\%$ ($p=4.81 \times 10^{-9}$) \\
5 & $-0.607$ & $2.16 \times 10^{-29}$ & $-0.043$ & 0.480 & $-4.18\%$ ($p=0.106$) \\
6 & $-0.540$ & $5.15 \times 10^{-22}$ & 0.071 & 0.239 & $+3.39\%$ ($p=0.097$) \\
7 & $-0.778$ & $1.44 \times 10^{-57}$ & 0.179 & $2.76 \times 10^{-3}$ & $+5.13\%$ ($p=9.70 \times 10^{-4}$) \\
8 & $-0.398$ & $2.13 \times 10^{-11}$ & 0.304 & $2.37 \times 10^{-7}$ & $+14.45\%$ ($p=4.88 \times 10^{-11}$) \\
9 & $-0.369$ & $2.06 \times 10^{-10}$ & 0.388 & $1.97 \times 10^{-11}$ & $+20.07\%$ ($p=8.15 \times 10^{-20}$) \\
10 & $-0.630$ & $4.28 \times 10^{-32}$ & 0.218 & $2.56 \times 10^{-4}$ & $+8.66\%$ ($p=1.43 \times 10^{-5}$) \\
\bottomrule
\end{tabular}
\end{table}

\begin{table}[p]
\centering
\caption{Female-only associations between static pelvic dimensions and size-scaled eigenvalues. Values are $R^2$ and $p$ for univariate linear models of $\log(\lambda_{\mathrm{scaled}})$ versus a single dimension.}
\label{tab:staticassoc}
\begin{tabular}{lcccc}
\toprule
& \multicolumn{2}{c}{Mode 9} & \multicolumn{2}{c}{Mode 2} \\
\cmidrule(lr){2-3}\cmidrule(lr){4-5}
Dimension & $R^2$ & $p$ & $R^2$ & $p$ \\
\midrule
PIT & 0.135 & $3.63 \times 10^{-6}$ & 0.089 & $2.11 \times 10^{-4}$ \\
AP & 0.092 & $1.59 \times 10^{-4}$ & 0.001 & 0.64 \\
Iliopectineal eminence width & 0.086 & $2.77 \times 10^{-4}$ & 0.026 & 0.050 \\
Sacral width & 0.077 & $5.96 \times 10^{-4}$ & 0.003 & 0.49 \\
Biacetabular width & 0.018 & 0.11 & 0.053 & $4.58 \times 10^{-3}$ \\
Biischiadic width & 0.012 & 0.19 & 0.060 & $2.55 \times 10^{-3}$ \\
Bituberous width & 0.013 & 0.17 & 0.078 & $5.57 \times 10^{-4}$ \\
Subpubic angle & 0.071 & $9.84 \times 10^{-4}$ & 0.015 & 0.14 \\
\bottomrule
\end{tabular}
\end{table}

\begin{table}[p]
\centering
\caption{Mean modal participation fractions (\%) by loading scenario after normalisation over modes~1--10.}
\label{tab:loadcases}
\begin{tabular}{llrrrrrrrrrr}
\toprule
Scenario & Sex & M1 & M2 & M3 & M4 & M5 & M6 & M7 & M8 & M9 & M10 \\
\midrule
Bilateral stance & Female & 0.4 & 84.6 & 0.1 & 7.3 & 4.0 & 1.0 & 0.1 & 1.3 & 1.1 & 0.1 \\
Bilateral stance & Male   & 0.6 & 84.5 & 0.1 & 4.9 & 4.2 & 1.6 & 0.2 & 2.9 & 0.8 & 0.1 \\
Unilateral stance & Female & 39.6 & 37.8 & 15.3 & 2.6 & 1.5 & 1.2 & 0.2 & 0.5 & 0.7 & 0.6 \\
Unilateral stance & Male   & 40.4 & 37.2 & 14.1 & 1.9 & 1.4 & 2.3 & 0.2 & 1.0 & 0.7 & 0.9 \\
Inlet opening & Female & 0.6 & 4.4 & 4.5 & 18.2 & 3.8 & 1.6 & 1.0 & 15.3 & 46.4 & 4.3 \\
Inlet opening & Male   & 0.7 & 3.1 & 5.1 & 6.0 & 3.4 & 1.5 & 1.1 & 18.1 & 52.5 & 8.5 \\
Outlet opening & Female & 0.4 & 2.6 & 0.4 & 15.7 & 40.7 & 17.8 & 0.9 & 3.8 & 16.1 & 1.6 \\
Outlet opening & Male   & 0.4 & 2.4 & 1.1 & 16.4 & 39.2 & 22.1 & 1.2 & 2.8 & 12.4 & 2.1 \\
Terminal descent & Female & 0.3 & 77.2 & 0.2 & 8.8 & 1.6 & 0.7 & 0.5 & 10.6 & 0.1 & 0.1 \\
Terminal descent & Male   & 0.5 & 74.7 & 0.2 & 4.8 & 1.8 & 0.7 & 0.7 & 16.4 & 0.1 & 0.1 \\
\bottomrule
\end{tabular}
\end{table}

\begin{table}[p]
\centering
\caption{Sensitivity hierarchy for mode~9, with reference values for modes~5 and~2. Values are cohort means.}
\label{tab:sensitivity}
\begin{tabular}{lccc}
\toprule
Structure & Mode 9 sensitivity & Mode 5 sensitivity & Mode 2 sensitivity \\
\midrule
Left SIJ cartilage interface & $3.32 \times 10^{-3}$ & $6.63 \times 10^{-4}$ & $2.24 \times 10^{-4}$ \\
Right SIJ cartilage interface & $2.69 \times 10^{-3}$ & $3.74 \times 10^{-4}$ & $1.97 \times 10^{-4}$ \\
Pubic symphysis interface & $2.07 \times 10^{-3}$ & $6.93 \times 10^{-4}$ & $9.45 \times 10^{-6}$ \\
Symphyseal ligament group & $1.09 \times 10^{-6}$ & $7.86 \times 10^{-8}$ & $2.12 \times 10^{-9}$ \\
Anterior SIJ ligaments (bilateral mean) & $4.64 \times 10^{-7}$ & $9.77 \times 10^{-8}$ & $2.59 \times 10^{-8}$ \\
Interosseous SIJ ligaments (bilateral mean) & $2.40 \times 10^{-7}$ & $4.31 \times 10^{-9}$ & $4.68 \times 10^{-11}$ \\
Posterior SIJ ligaments (bilateral mean) & $1.27 \times 10^{-8}$ & $3.47 \times 10^{-8}$ & $2.87 \times 10^{-9}$ \\
Sacrospinous ligaments & $6.57 \times 10^{-9}$ & $3.07 \times 10^{-9}$ & $5.85 \times 10^{-9}$ \\
Sacrotuberous ligaments & $2.37 \times 10^{-8}$ & $1.24 \times 10^{-9}$ & $6.49 \times 10^{-10}$ \\
\bottomrule
\end{tabular}
\end{table}

\begin{figure}[p]
\centering
\fbox{\parbox[c][7cm][c]{0.9\linewidth}{\centering Placeholder for Figure~1\\
Cohort-to-template workflow and correspondence-preserving finite-element pipeline.}}
\caption{Study workflow. Clinical CT-derived pelvic anatomies were mapped onto a common template, converted into subject-specific finite-element models, decomposed into elastic eigenmodes, and analysed through load-case projection and sensitivity analysis. The key output is not a single optimal pelvic shape, but a population-scale mechanical phenotype defined by accessible deformation pathways.}
\label{fig:workflow}
\end{figure}

\begin{figure}[p]
\centering
\fbox{\parbox[c][7cm][c]{0.9\linewidth}{\centering Placeholder for Figure~2\\
Size-corrected modal spectrum and adjusted sex effects.}}
\caption{Size-corrected modal spectrum. Modes~2 and~9 show the largest adjusted sex effects after control for pelvic scale and age, whereas mode~1 acts as a near-neutral reference. This pattern indicates that female compliance is not merely a size effect but part of a more specific reorganisation of ring mechanics.}
\label{fig:eigs}
\end{figure}

\begin{figure}[p]
\centering
\fbox{\parbox[c][8cm][c]{0.9\linewidth}{\centering Placeholder for Figure~3\\
Representative deformation maps of modes~1, 2, 5 and~9.}}
\caption{Representative low-order deformation pathways. Mode~1 captures a conserved unilateral support component, mode~2 the baseline support-related ring mode, mode~5 an outlet-related pathway, and mode~9 the principal candidate inlet-opening pathway. The biological interpretation arises from this modal hierarchy rather than from any single dimension alone.}
\label{fig:modeshapes}
\end{figure}

\begin{figure}[p]
\centering
\fbox{\parbox[c][8cm][c]{0.9\linewidth}{\centering Placeholder for Figure~4\\
Stacked-bar modal participation fractions by loading scenario.}}
\caption{Load-case modal recruitment. Habitual support scenarios are built on a shared low-order modal backbone in both sexes. Under simulated inlet opening, men concentrate more of the response in mode~9, whereas women distribute the response more broadly across coordinated low-order pathways, consistent with a distributed compliance strategy.}
\label{fig:recruitment}
\end{figure}

\begin{figure}[p]
\centering
\fbox{\parbox[c][7cm][c]{0.9\linewidth}{\centering Placeholder for Figure~5\\
Mode~9 versus static pelvic dimensions in women.}}
\caption{Static pelvimetry captures only a modest part of inlet-opening mechanics. PIT and AP are the strongest single predictors of size-corrected mode~9 stiffness in women, yet no single dimension explains more than a modest fraction of the variance.}
\label{fig:pelvimetry}
\end{figure}

\begin{figure}[p]
\centering
\fbox{\parbox[c][7cm][c]{0.9\linewidth}{\centering Placeholder for Figure~6\\
Sensitivity hierarchy of mode~9 across interfaces and ligament groups.}}
\caption{Sensitivity hierarchy for the principal inlet-opening pathway. Mode~9 is governed primarily by the pubic symphysis and anteriorly weighted sacroiliac constraints, whereas posterior ligament groups are much less influential. The inlet-opening pathway is therefore best understood as an anteriorly regulated ring mechanism rather than as diffuse pelvic laxity.}
\label{fig:sensitivity}
\end{figure}

\clearpage
\printbibliography

\end{document}
